% Objections and replies for "The Unchosen Adjacency"

\section{Objections and Replies}\label{sec:objections}

Five objections are considered. Each is stated in its strongest form. Two of the replies concede something, and the concessions are carried into the framework and not merely acknowledged.

\subsection{``A private platform compels nothing; users can leave''}

The objection is correct about doctrine. Constitutional compelled association requires state action, and a user who joins a private service on its terms is not compelled in that sense. The essay does not claim otherwise, which is why it uses \emph{unchosen bundling} as its general term.

The objection does not reach the design claim. That a user may exit does not show that the interface has offered the separations the user needed. Under the exit proposition, when an intention is withheld, every available action carries consequential mismatch with it, and the only way to avoid all of it is to forgo the change. Exit is then the cost the interface imposes, and its availability describes the problem without answering it. Nothing here depends on the exit being difficult. The claim holds even where leaving is easy.

\subsection{``Advertising is the business model, and disaggregation costs revenue''}

The objection is partly right, and the framework should say so. The defective-action test asks whether a cleaner action is technically feasible, not whether providing it is costless. A platform may have sound commercial reasons for withholding an action.

Two replies apply. First, feasibility is not merely asserted. Where a comparator exists, including advertiser-side controls over adjacency, the capability has been built and deployed. Second, a commercial justification is a reason for the design, and the test classifies the design without denying the reason. The framework therefore \emph{classifies} bundles and does not by itself \emph{mandate} their removal. Whether a defective action should be corrected is a further question that weighs the user-side cost of the bundle against the platform-side cost of the separation. The essay reports the first term and does not claim to have measured the second.

\subsection{``Simple interfaces are a virtue''}

A single block button is easy to understand. Each added control adds options, documentation, and error, and technical feasibility does not oblige a designer to expose every action.

The objection identifies a real cost that the definitions omit: $\mathcal{A}^{*}$ counts what can be offered, not what can be offered without degrading usability. The reply is that simplicity and disaggregation are not opposed by nature. Bundling is a simplification that hides a choice, and the hidden choice is the one the user needs to make. Defaults, presets, and graduated escalation are familiar ways to reduce the burden of finer controls, so that the finer action is available without being demanded, although the essay does not survey them. The essay concedes, however, that the counterfactual criterion identifies where a better interface is possible and does not by itself show that it is better overall.

\subsection{``Separating perception from participation enables stalking and brigading''}

The strongest form is that a person barred from commenting but able to watch can still harass through other means, such as a second account, screenshots, or coordination elsewhere. Blocking exists precisely because visibility can itself be a harm.

The framework accommodates this and does not hide it. The loss criterion requires that the user place positive value on the retained relation. When a user wants the other person to lose sight of them, removing perception belongs to the intended set $S_a$ and is not surplus. A block used against stalking or threats is therefore not a defective action, and the essay says that block remains the correct action in those cases. The circumvention objection also applies equally to blocks, which second accounts likewise evade, so it does not favor the coarser control. What the framework requires is that the two purposes be distinguishable, so that the user who wants only to bar participation is not forced into the remedy suited to the user who needs to disappear.

\subsection{``The thresholds are arbitrary, so any verdict can be produced''}

The parameters $\theta_\kappa$, $\rho_\kappa$, and $\lambda_t$ are normative and unestimated, and a framework that returns whatever verdict its author wants would be worthless.

Three replies apply. First, the parameters are declared and typed by claim, so a verdict can be traced to a stated choice. Second, the framework does not claim a verdict for every case. The applications section returns different results for different cases and excludes one, which is not the behavior of an instrument that condemns everything. Third, the appropriate test of the result is whether the ordering of cases survives changes in the parameters. The applications section states a condition: if attribution for adjacency stochastically dominates attribution for presence, and the thresholds for institutional approval are at least as demanding as those for commercial endorsement, then the ordering cannot reverse. Whether those assumptions hold is an empirical and normative question that the essay does not settle. If they fail, the ordering may flip, and the essay should then report that.

\subsection{What the replies change}

Two concessions carry forward. The framework classifies mismatches and does not itself mandate their removal, which affects how the design implications are framed. And the counterfactual criterion identifies possible improvement without establishing that the improvement is net beneficial. The safety objection is handled inside the formalism, not by exception, through the requirement that the user value the relation that is lost.
